针对银行需要在有信贷记录企业、无信贷记录企业及突发压力情景下制定授信策略的问题，本文构建“经营特征—信用风险—利率响应—组合决策”的递进框架。基于企业信息和进销项发票，在统一时间窗口内提取经营特征；以人工信誉评级为监督标签，使用 Logistic 与 XGBoost 的折外概率融合构造信用风险指数（CRI）。

在无信贷记录企业上，本文仅使用两批企业均可获得的发票经营特征进行共同特征跨样本外推，并以漂移诊断和人工复核控制应用风险；不将该步骤表述为领域自适应或校准后的违约概率模型。随后，利用保序回归描述利率与客户留存的关系，并将风险、流失、额度与预算约束纳入组合授信配置。

在有限压力情景下，模型以最坏情景收益为目标进行稳健配置。该部分定位为情景压力测试与有限情景最大最小决策，不宣称连续不确定集或分布鲁棒优化理论创新。本文不把评级概率解释为经校准的违约概率；模型的定位是提供可解释、可复核的风险排序和授信配置依据，为后续人工复核和授信管理提供可复核的决策依据。
